\worldchapter{Fortschritt und Auszeit}{advancement}
\label{ch:fortschritt}

\begin{multicols*}{2}
\section{Grundsatz}
Fortschritt\index{Fortschritt} belohnt überstandene Kosten, Verwundbarkeit und Reputation, nicht stumpfes Wiederholen.

\section{Wachstum}
Am Ende einer Sitzung erhält ein Charakter 1 \textbf{Wachstum}\index{Growth}, wenn mindestens einer der folgenden Punkte zutraf:
\begin{itemize}[leftmargin=*]
\item sie hat eine große Komplikation aus einer Verwundbarkeit akzeptiert
\item sie hat eine \textbf{kritische Konsequenz} überlebt, die mit einem Mark verknüpft war
\end{itemize}

Bei 3 Wachstum wählst du genau eines:
\begin{itemize}[leftmargin=*]
\item eine Fertigkeit von 0 auf 1 erhöhen
\item eine Fertigkeit von 1 auf 2 erhöhen
\item eine Fertigkeit von 2 auf 3 erhöhen
\item eine \textbf{Wegfacette} ergänzen
\item 1 Verderbnis entfernen und 1 neue Verpflichtung erhalten
\end{itemize}

\section{Auszeit}
Zwischen großen Einsätzen erhält jeder Charakter \textbf{2 Auszeit-Aktionen}\index{Auszeit}.
Wiederholst du dieselbe Aktion in derselben Auszeit, bekommt jeder Wiederholungsversuch kumulativ -10.

\section{Auszeit-Aktionen}

Während der Auszeit kann ein Charakter eine der folgenden Aktionen ausführen:

\begin{itemize}[leftmargin=*]
\item \textbf{Bindung pflegen}
\item \textbf{Wunden versorgen}
\item \textbf{Kraft schöpfen}
\item \textbf{Gunst suchen}
\item \textbf{Nachforschungen anstellen}
\item \textbf{Arkana sammeln}
\item \textbf{Reinigung suchen}
\item \textbf{Ausrüstung beschaffen}
\end{itemize}

\subsection{Erholen}\index{Erholen}
\[
Wille + Körper
\]
Erfolg: Entferne 2 Wunden oder 3 Bürde oder 1 Wunde + 1 Bürde + 1 Zustand.\\
Misserfolg: Entferne nur 1 Bürde.

\subsection{Reparieren}\index{Reparieren}
\[
Geist oder Geschick + Handwerk
\]

Erfolg: Repariere einen zerstörten Gegenstand.\\
Starker oder kritischer Erfolg: Repariere stattdessen zwei Gegenstände.\\
Misserfolg: Die Reparatur scheitert. Erst neues Material, Bezahlung oder zusätzliche Zeit erlauben einen weiteren Versuch.

\subsection{Acquire}\index{Acquire}
\[
Mind oder Appearance + Influence
\]
Erfolg: ein nützlicher Gegenstand, Dienst oder zeitweiliger Kontakt.\\
Starker oder kritischer Erfolg: zwei Dinge oder ein hochwertiger Vermögensvorteil.\\
Misserfolg: Preis verdoppelt sich, Gefallen wird geschuldet oder Misstrauen wächst.

\subsection{Beschaffen}\index{Beschaffen}
\[
Geist oder Erscheinung + Einfluss
\]

Erfolg: Beschaffe einen nützlichen Gegenstand, einen Dienst oder einen zeitweiligen Kontakt.\\
Starker oder kritischer Erfolg: Beschaffe stattdessen zwei solcher Dinge oder einen bedeutenden Vermögensvorteil.\\
Misserfolg: Der Preis verdoppelt sich, ein Gefallen wird geschuldet oder Misstrauen entsteht.

\subsection{Trainieren}\index{Traineren}
\[
Wille + Resolve
\]
Erfolg: 1 Trainingsvermerk notieren.\\
Bei 3 Trainingsvermerken: Ein Attribut um 1 steigern, jedoch nicht über 4.\\
Misserfolg: kein Vermerk; bei kritischem Misserfolg +1 Burden.

\subsection{Kontakte pflegen}\index{Kontakte pflegen}
\[
Erscheinung + Einfluss \text{ oder } Wille + Standhaftigkeit
\]

Erfolg: Erneuere eine belastete Bindung oder sichere dir einen konkreten Gefallen.\\
Starker oder kritischer Erfolg: Beides.\\
Misserfolg: Der Kontakt fordert einen Preis, eine Verpflichtung oder zusätzliche Zeit.

\subsection{Ritualvorbereitung}\index{Ritualvorbereitung}
\[
Geist oder Wille + Gabe oder Glaubensniveau
\]
Erfolg: ein Ritualvorteil, +10 auf einen zukünftigen passenden Zauber oder eine Anrufung.\\
Starker oder kritischer Erfolg: zusätzlich 1 Arkana oder 1 Gunst regenerieren.\\
Misserfolg: +1 Bürde.

\section{Verpflichtungen}
Eine \textbf{Verpflichtung}\index{Verpflichtung} ist eine dauerhafte Schuld, Pflicht, ein Versprechen oder ein Dienst.
\begin{itemize}[leftmargin=*]
\item einmal pro Sitzung darf die Spielleiterin sie einfordern
\item akzeptieren: +1 Omen und die Komplikation tragen
\item verweigern: +1 Burden und die zugehörige Beziehung verschlechtert sich
\end{itemize}

\section{Auffüllen von Ressourcen}
Am Ende der Auszeit:
\begin{itemize}[leftmargin=*]
\item Omen auf Max auffüllen
\item Arkana Pool und Gunst voll auffüllen
\end{itemize}
\end{multicols*}
